\section{Prediction-Error Curiosity}
\label{sec:prediction_error_curiosity}

Prediction-error curiosity rewards transitions that are difficult for a learned model to predict.
The general assumption is that high prediction error indicates novelty, surprise, or incomplete understanding of the environment.
If a model predicts the next representation $\hat{z}_{t+1}=f_\psi(z_t,a_t)$ from the current representation $z_t$ and action $a_t$, an intrinsic reward can be defined as
\begin{equation}
    r_t^{\mathrm{int}} = \left\lVert f_\psi(z_t,a_t)-z_{t+1}\right\rVert_2^2.
\end{equation}
Early deep predictive approaches used this idea to encourage agents to seek transitions that improve their understanding of environment dynamics \parencite{stadie-2015}.

The Intrinsic Curiosity Module (ICM) is a representative feature-space prediction-error method \parencite{pathak-2017}.
ICM learns an encoder $\phi_\omega$ that maps observations to features $z_t=\phi_\omega(o_t)$, a forward model that predicts $z_{t+1}$ from $(z_t,a_t)$, and an inverse model that predicts $a_t$ from $(z_t,z_{t+1})$.
The inverse model encourages the representation to retain information about controllable aspects of the transition, reducing sensitivity to observation details that do not depend on the agent's action.
The curiosity reward is then based on forward prediction error in this learned feature space rather than raw observation space.

Random Network Distillation (RND) provides another influential form of prediction-error curiosity \parencite{burda-2018B}.
RND uses a fixed random target network $h$ and a trainable predictor $\hat{h}_\psi$.
The intrinsic reward is the prediction error
\begin{equation}
    r_t^{\mathrm{int}} = \left\lVert \hat{h}_\psi(o_t)-h(o_t)\right\rVert_2^2.
\end{equation}
Because the target network is fixed, prediction error tends to be high for unfamiliar observations and decreases as the predictor trains on visited states.
Large-scale studies of curiosity-driven learning have shown that such prediction-error signals can support exploration across a range of environments \parencite{burda-2018A}.

Prediction-error curiosity is powerful because it does not require explicit state counting and can be implemented with neural networks.
It can discover sparse rewards by making unfamiliar observations attractive, and it can be combined with standard policy-gradient methods through reward augmentation.
However, its main weakness is that prediction error conflates learnable novelty with irreducible uncertainty.
A transition may remain difficult to predict because it is stochastic, partially observed, chaotic, or unrelated to the agent's controllable behavior.
In such cases, prediction error can remain high even after extensive experience.

This weakness is closely related to the noisy-TV problem.
If an agent can repeatedly trigger observations that are unpredictable but behaviorally unproductive, a prediction-error bonus may reward this behavior indefinitely.
Feature learning and inverse dynamics can reduce the problem by emphasizing controllable factors, but they do not eliminate all sources of unlearnable error.
RND can also be affected by observation streams that maintain high predictor error for reasons unrelated to task progress.

Prediction-error methods therefore motivate the use of temporal improvement as a refinement of curiosity.
High error alone answers the question of what the model currently fails to predict; it does not answer whether additional data are making the model better.
Learning-progress methods address this distinction by rewarding reductions in error rather than the absolute error value.